\vspace{-5mm}
\section{Related Work}

\label{sec:relatedWork}
%In this paper, we focus on the problem of implementing effective invariant enforcement for attack prevention in smart contracts.
% The related works can be broadly classified into security analysis and runtime validation of smart
% contracts.





\noindent
\textbf{Smart Contract Invariants.} 
Prior works have studied various smart contract invariants. \citeauthor{zhou2020ever}~\cite{zhou2020ever} introduces six invariants to defend against different hacks. Cider~\cite{liu2022learning} uses deep reinforcement learning to learn invariants that prevent arithmetic overflows, while SPCon~\cite{liu2022finding} recovers likely access control models from function callers in past transactions. Over~\cite{deng2024safeguarding} infers safety constraints on oracles using contract source code and oracle update history. In contrast, {\name} explores a broader range of invariants and assesses their effectiveness and bypassability against real-world attacks.

% %Prior works have delved into various types without always explicitly using the term "invariant". 
% Prior works~\cite{liu2022learning,liu2022finding,deng2024safeguarding,zhou2020ever} studies specific types of smart contract invariants.
% \citeauthor{zhou2020ever}~\cite{zhou2020ever} introduces $6$ invariants to defend against 
% different hacks. Cider~\cite{liu2022learning} leverages deep reinforcement learning on smart contract source code 
% to learn invariants that prevent arithmetic overflows. SPCon~\cite{liu2022finding} utilizes function callers in past transactions of a contract to recover  
% a likely access control model. More recently, Over~\cite{deng2024safeguarding} infers safety constraints on oracles from the contract's source code and the history of oracle updates.
% Compared to the above works, {\name} studies a broader range of 
% invariants and evaluates their effectiveness and bypassability 
% against real-world attacks.}



% SPCon~\cite{liu2022finding} utilizes function callers in past transactions of a contract to remover  
% a likely access control model. These two works have shown the invariant effectiveness on their 
% respective categories. However, they do not cover the full spectrum of invariants as we do. 
% InvCon~\cite{liu2022invcon} and its followup work, InvCon+~\cite{liu2024automated}, 
% closely align with our methodology by also deriving many types of invariants 
% from past transactions. 
% They leverage pre/post conditions of transactions to generate all possible invariants. 
% However, it results in thousands of invariants for each contract, which is impractical for real-world deployment,
% as illustrated in our comparative study in RQ5 in Section~\ref{sec:evaluation}.
% }



\noindent
\textbf{Smart Contract Security Analysis.}
%There are a large body of works on vulnerabiltiy detection of smart contracts~\cite{bibid}.
%They have demonstrated powerful ability to find common smart contract vulnerabilities, e.g., integer overflow/underflow~\cite{BECoverflow}, reentrancy~\cite{DAOattacks}, dangerous delegatecall~\cite{Paritystolen}, and etc.
%These vulnerabilities are usually critical for smart contract security.
%For instance, the Parity wallet was hacked in 2017 due to the missing protection of delegatecall operation.
%As a result, the attacker took over the ownership of the wallet and then stole 150, 000 ETH worth 30 million USD.
There is a large body of works on the detection of security
vulnerabilities in smart contracts.
%They have demonstrated powerful ability to find common smart contract vulnerabilities.
Oyente~\cite{luu2016making} is one of the first symbolic execution-based security tools to detect \textit{reentrancy}, \textit{mishandled exception}, \textit{transaction order dependency}, and \textit{timestamp dependency}.
It has been extended to detect \textit{greedy}, \textit{prodigal} and \textit{suicidal}
contracts~\cite{nikolic2018finding}.
Other well-known symbolic-execution tools include Manticore~\cite{manticore} and
Mythril~\cite{mythril} which are able to find other types of vulnerabilities, e.g.,
\textit{dangerous delegatecall}, \textit{integer overflow}, etc.
Slither~\cite{slither} is another popular static security analysis tool for smart contracts.
It performs data flow and control flow dependency analysis to support up to $87$ bug detectors.
Other static analyzers include SmartCheck~\cite{tikhomirov2018smartcheck} mainly targeting
\textit{bad coding practices}, Securify~\cite{securify} and Ethainter~\cite{brent2020ethainter}
for finding \textit{information-flow} vulnerabilities.
Moreover, dynamic analysis tools~\cite{2018contractfuzzer, atzei2017survey, nguyen2020sfuzz,
echidna, wustholz2020harvey, wang2020oracle, liu2020modcon} were proposed to detect smart contract
vulnerabilities through fuzzing and model-based testing.
%ContractFuzzer~\cite{2018contractfuzzer} is one of the earliest black-box fuzzing tools for smart contracts to detect seven common vulnerabilities listed in~\cite{atzei2017survey}.
%It defines practical test oracles and analyzes the collected execution traces to check against the oracles.
%The gray-box fuzzing tools include sFuzz~\cite{nguyen2020sfuzz}, Echinda~\cite{echidna}, Harvey~\cite{wustholz2020harvey}, and ContraMaster~\cite{wang2020oracle}.
%All of these tools employ some kinds of feed-back mechanisms to guide their testing process.
%\textcolor{green}{Liu Ye:
%The techniques used to address the security issues can be classified into \emph{static} and
%\emph{dynamic} analyses.
%The former can be further broken down to program analysis, program synthesis, symbolic execution,
%formal verification, and theorem proving.
%Slither~\cite{feist2019slither} and Ethainter~\cite{brent2020ethainter} perform taint analysis to
%find information flow vulnerabilities.
%SmartCheck~\cite{tikhomirov2018smartcheck} targets code issues by searching contracts' AST against
%predefined rules.
%Securify~\cite{securify} infers semantic information by analyzing the control- and
%data-dependencies of contract code, which is then checked against several predefined security
%patterns.
%Meanwhile, SmartScopy~\cite{feng2019precise} introduces a summary-based symbolic evaluation to
%synthesize attack programs for vulnerable contracts.
%Oyente~\cite{luu2016making} is one of the earliest symbolic execution engine for smart contracts,
%followed by Manticore~\cite{manticore} and Mythril~\cite{mythril}.
%Mythril~\cite{mythril} is an industrial security analysis tool which combines symbolic execution and
%taint analysis to detect nearly 30 classes of vulnerabilities.
%Other formal verification~\cite{wang2018formal,kalra2018zeus,liu2020towards} and
%theorem-proving~\cite{hirai2017defining,permenev2020verx} tools were
%built aiming to check properties on the safety~\cite{kalra2018zeus,permenev2020verx}, security, and
%fairness~\cite{liu2020towards} of smart contracts.}
%
%\textcolor{green}{Liu Ye: On the other hand, dynamic analyses find exploitable security bugs by directly executing the
%contracts in a testing environment.
%Test cases are usually generated through fuzzing or model-based testing~\cite{liu2020modcon}.
%ContractFuzzer~\cite{2018contractfuzzer} is one of the earliest black-box fuzzing frameworks
%for detecting smart contract vulnerabilities.
%It predefines several practical test oracles and analyzes the collected execution traces to check
%against the oracles.
%The gray-box fuzzing tools ContraMaster~\cite{wang2020oracle,wang2019vultron},
%Harvey~\cite{wustholz2020harvey}, and Echidna~\cite{echidna} employ a feedback-driven mechanism to
%guide the testing process.
%}
%Especially, with the advent of dencentralized finance application (DeFi), billions of dollars worth asset is run by these applications.
%Although, according to the hack statistics~\cite{defillama}, as of 12 September 2023, DeFi hacks brought 5.53 billion dollars out of the overall 6.94 billion dollars loss in block incidents.
%In addition, DeFi is an interdisciplinary domain that has its own attack surface~\cite{zhou2023sok}.
%Different from the abovementioned common vulnerabilities,
%DeFi vulnerabilities are more subtle and many hacks are designed in a much more sophisticated manner.
%Serious DeFi attack types include governance attack, oracle price manipulation, and front-running~\cite{werner2022sok, zhou2023sok}.
%To achieve decentralization, DeFi protocols tend to implement decentralized governance mechanisms based on governance tokens used for proposing and voting on protocol upgrades.
%However, once attackers had sufficient amount of governance tokens to execute malicious proposals,
%the security of the whole DeFi protocols is likely undermined.
%Oracle price manipulation is another common attack type.
%To provide service of asset trading, every DeFi application relies on price oracles, mostly provided by AMMs (Automated Market Makers) for decentralization purpose.
%Within oracle price manipulation,
%attackers often first borrow a large amount of funds through flashloan mechanism~\cite{bibid} and then use part of these funds to create an imbalance in the liquidity pools within an AMM, leading to the price slippage.
%Since the trading within DeFi applications depends on AMM's price,
%attackers can exploit the created price gap to make profit through buying DeFi's assets and repaying the flashloan.
%Front-running within blockchain is such kind of attack where a transaction is submitted and purposes to be executed before some pending user transactions.
%Alike technique also includes back-running~\cite{bibid} and the sandwich attacks~\cite{bibid} can be ascribed to them.

There is a substantial body of work focusing on the detection and exploitation of DeFi vulnerabilities. \citeauthor{qin2021attacking}~\cite{qin2021attacking} and \citeauthor{zhou2021just}~\cite{zhou2021just} manually formulated vulnerabilities such as \textit{oracle price manipulation} and \textit{arbitrage} into an optimization problem. \camera{FlashSyn~\cite{chen2022flashsyn} proposed a framework that automatically formulates flash loan attacks as a synthesis problem, by approximating the functions of DeFi protocols.}
\citeauthor{wu2021defiranger}~\cite{wu2021defiranger} identified several \textit{oracle price
manipulation} patterns from on-chain transaction data to detect real-world attacks while
\citeauthor{kong2023defitainter}~\cite{kong2023defitainter} detects price manipulation
vulnerabilities in DeFi applications through inter-contract taint analysis.
Also, \citeauthor{gudgeon2020decentralized}~\cite{gudgeon2020decentralized} showcased how to
use \textit{flashloan} to conduct \textit{governance attack}.
\citeauthor{baum2022sok}~\cite{baum2022sok} surveyed the state-of-the-art mitigation techniques for
\textit{front-running} in DeFi.
Interestingly, several works~\cite{xue2022preventing, deng2023robust, zhang2023your} explored
front-running as a defense mechanism against smart contract exploits.
Different from the over-generalized security patterns used by the existing tools,
our invariant guards capture the subtle semantic constraints of specific smart contracts.

%Most of the security analysis works are pattern-based and focus on vulnerability detection.
%Their effectiveness are subjected to the simple and limited patterns which are not able to capture sophisticated attacks.
%In this paper, we automatically synthesize the normal behavior patterns from benign transaction histories through dynamic invariant inference and aim to apply these learned patterns to filter out abnormal transactions which are potentially the attacks.


\noindent
\textbf{Runtime Verification and Validation.} Runtime verification has been used for validation purpose where the runtime checks are on properties 
from users' expectation rather than from formal program specifications~\cite{magazzeni2017validation}.
Sereum~\cite{rodler2018sereum} is a general runtime validation framework to protect 
deployed contracts against reentrancy attacks.
Solythesis~\cite{li2020securing} provides a source-to-source compiler that compiles smart contracts with user-specified invariants to reject unexpected transactions~\cite{li2020securing}.
In contrast, {\name} is not limited to predefined attack types and has demonstrated effectiveness in mitigating a wide range of sophisticated attacks.





% \subsection{Dynamic Invariant Detection}
% Invariants are program properties that always hold during program execution.
% Static and dynamic invariant detection are two fundamental approaches used in program analysis to
% identify program invariants.
% Static invariant detection involves analyzing the source code of a program without executing it.
% %The goal is to identify invariants that hold under all possible program executions. Techniques such as abstract interpretation, static analysis, and theorem proving are commonly employed in this context. Static invariant detection is useful for reasoning about program behavior in a broader context and can help in bug detection, optimization, and program verification.
% Due to the inherent complexity of analyzing all possible program paths statically, this approach
% may struggle with modeling certain complex program features and may produce imprecise results.

% Dynamic invariant detection, on the other hand, relies on observing the behavior of a program
% during its execution.
% It involves monitoring the program's states and tracking its behaviors to infer invariants that
% hold for specific execution instances.
% %Techniques like runtime assertion checking, program slicing, and dataflow analysis are often used in dynamic invariant detection.
% This approach provides more accurate and context-specific invariants, as it has access to the
% actual runtime behaviors of the program.
% Daikon~\cite{ernst2007daikon}, a well-known dynamic invariant detection tool, is able to detect
% likely program invariants for Java programs from their test case executions.
% Daikon takes as input a set of collected program data traces and then checks them against the
% predefined invariant templates.
% The invariant candidates are discarded if they are refuted by any of the data traces.
% The resulting \emph{likely} invariants have been proved valuable for debugging, testing, and
% runtime verification of programs~\cite{bibid}.
% InvCon~\cite{liu2022invcon} extends Daikon by supporting dynamic invariant detection for Solidity
% smart contracts, based on their past transaction histories.
% %However, it incurs some overhead due to the runtime monitoring of the whole contract state.
% %Dynamic invariant detection is particularly valuable for debugging, testing, and runtime verification of programs.


